\IfFileExists{stacks-project.cls}{%
\documentclass{stacks-project}
}{%
\documentclass{amsart}
}

% For dealing with references we use the comment environment
\usepackage{verbatim}
\newenvironment{reference}{\comment}{\endcomment}
%\newenvironment{reference}{}{}
\newenvironment{slogan}{\comment}{\endcomment}
\newenvironment{history}{\comment}{\endcomment}

% For commutative diagrams we use Xy-pic
\usepackage[all]{xy}

% We use 2cell for 2-commutative diagrams.
\xyoption{2cell}
\UseAllTwocells

% We use multicol for the list of chapters between chapters
\usepackage{multicol}

% This is generally recommended for better output
\usepackage{lmodern}
\usepackage[T1]{fontenc}

% For cross-file-references
\usepackage{xr-hyper}

% Package for hypertext links:
\usepackage{hyperref}

% For any local file, say "hello.tex" you want to link to please
% use \externaldocument[hello-]{hello}
\externaldocument[introduction-]{introduction}
\externaldocument[conventions-]{conventions}
\externaldocument[sets-]{sets}
\externaldocument[categories-]{categories}
\externaldocument[topology-]{topology}
\externaldocument[sheaves-]{sheaves}
\externaldocument[sites-]{sites}
\externaldocument[stacks-]{stacks}
\externaldocument[fields-]{fields}
\externaldocument[algebra-]{algebra}
\externaldocument[brauer-]{brauer}
\externaldocument[homology-]{homology}
\externaldocument[derived-]{derived}
\externaldocument[simplicial-]{simplicial}
\externaldocument[more-algebra-]{more-algebra}
\externaldocument[smoothing-]{smoothing}
\externaldocument[modules-]{modules}
\externaldocument[sites-modules-]{sites-modules}
\externaldocument[injectives-]{injectives}
\externaldocument[cohomology-]{cohomology}
\externaldocument[sites-cohomology-]{sites-cohomology}
\externaldocument[dga-]{dga}
\externaldocument[dpa-]{dpa}
\externaldocument[sdga-]{sdga}
\externaldocument[hypercovering-]{hypercovering}
\externaldocument[schemes-]{schemes}
\externaldocument[constructions-]{constructions}
\externaldocument[properties-]{properties}
\externaldocument[morphisms-]{morphisms}
\externaldocument[coherent-]{coherent}
\externaldocument[divisors-]{divisors}
\externaldocument[limits-]{limits}
\externaldocument[varieties-]{varieties}
\externaldocument[topologies-]{topologies}
\externaldocument[descent-]{descent}
\externaldocument[perfect-]{perfect}
\externaldocument[more-morphisms-]{more-morphisms}
\externaldocument[flat-]{flat}
\externaldocument[groupoids-]{groupoids}
\externaldocument[more-groupoids-]{more-groupoids}
\externaldocument[etale-]{etale}
\externaldocument[chow-]{chow}
\externaldocument[intersection-]{intersection}
\externaldocument[pic-]{pic}
\externaldocument[weil-]{weil}
\externaldocument[adequate-]{adequate}
\externaldocument[dualizing-]{dualizing}
\externaldocument[duality-]{duality}
\externaldocument[discriminant-]{discriminant}
\externaldocument[derham-]{derham}
\externaldocument[local-cohomology-]{local-cohomology}
\externaldocument[algebraization-]{algebraization}
\externaldocument[curves-]{curves}
\externaldocument[resolve-]{resolve}
\externaldocument[models-]{models}
\externaldocument[functors-]{functors}
\externaldocument[equiv-]{equiv}
\externaldocument[pione-]{pione}
\externaldocument[etale-cohomology-]{etale-cohomology}
\externaldocument[proetale-]{proetale}
\externaldocument[relative-cycles-]{relative-cycles}
\externaldocument[more-etale-]{more-etale}
\externaldocument[trace-]{trace}
\externaldocument[crystalline-]{crystalline}
\externaldocument[spaces-]{spaces}
\externaldocument[spaces-properties-]{spaces-properties}
\externaldocument[spaces-morphisms-]{spaces-morphisms}
\externaldocument[decent-spaces-]{decent-spaces}
\externaldocument[spaces-cohomology-]{spaces-cohomology}
\externaldocument[spaces-limits-]{spaces-limits}
\externaldocument[spaces-divisors-]{spaces-divisors}
\externaldocument[spaces-over-fields-]{spaces-over-fields}
\externaldocument[spaces-topologies-]{spaces-topologies}
\externaldocument[spaces-descent-]{spaces-descent}
\externaldocument[spaces-perfect-]{spaces-perfect}
\externaldocument[spaces-more-morphisms-]{spaces-more-morphisms}
\externaldocument[spaces-flat-]{spaces-flat}
\externaldocument[spaces-groupoids-]{spaces-groupoids}
\externaldocument[spaces-more-groupoids-]{spaces-more-groupoids}
\externaldocument[bootstrap-]{bootstrap}
\externaldocument[spaces-pushouts-]{spaces-pushouts}
\externaldocument[spaces-chow-]{spaces-chow}
\externaldocument[groupoids-quotients-]{groupoids-quotients}
\externaldocument[spaces-more-cohomology-]{spaces-more-cohomology}
\externaldocument[spaces-simplicial-]{spaces-simplicial}
\externaldocument[spaces-duality-]{spaces-duality}
\externaldocument[formal-spaces-]{formal-spaces}
\externaldocument[restricted-]{restricted}
\externaldocument[spaces-resolve-]{spaces-resolve}
\externaldocument[formal-defos-]{formal-defos}
\externaldocument[defos-]{defos}
\externaldocument[cotangent-]{cotangent}
\externaldocument[examples-defos-]{examples-defos}
\externaldocument[algebraic-]{algebraic}
\externaldocument[examples-stacks-]{examples-stacks}
\externaldocument[stacks-sheaves-]{stacks-sheaves}
\externaldocument[criteria-]{criteria}
\externaldocument[artin-]{artin}
\externaldocument[quot-]{quot}
\externaldocument[stacks-properties-]{stacks-properties}
\externaldocument[stacks-morphisms-]{stacks-morphisms}
\externaldocument[stacks-limits-]{stacks-limits}
\externaldocument[stacks-cohomology-]{stacks-cohomology}
\externaldocument[stacks-perfect-]{stacks-perfect}
\externaldocument[stacks-introduction-]{stacks-introduction}
\externaldocument[stacks-more-morphisms-]{stacks-more-morphisms}
\externaldocument[stacks-geometry-]{stacks-geometry}
\externaldocument[moduli-]{moduli}
\externaldocument[moduli-curves-]{moduli-curves}
\externaldocument[examples-]{examples}
\externaldocument[exercises-]{exercises}
\externaldocument[guide-]{guide}
\externaldocument[desirables-]{desirables}
\externaldocument[coding-]{coding}
\externaldocument[obsolete-]{obsolete}
\externaldocument[fdl-]{fdl}
\externaldocument[index-]{index}

% Theorem environments.
%
\theoremstyle{plain}
\newtheorem{theorem}[subsection]{Theorem}
\newtheorem{proposition}[subsection]{Proposition}
\newtheorem{lemma}[subsection]{Lemma}

\theoremstyle{definition}
\newtheorem{definition}[subsection]{Definition}
\newtheorem{example}[subsection]{Example}
\newtheorem{exercise}[subsection]{Exercise}
\newtheorem{situation}[subsection]{Situation}

\theoremstyle{remark}
\newtheorem{remark}[subsection]{Remark}
\newtheorem{remarks}[subsection]{Remarks}

\numberwithin{equation}{subsection}

% Macros
%
\def\lim{\mathop{\mathrm{lim}}\nolimits}
\def\colim{\mathop{\mathrm{colim}}\nolimits}
\def\Spec{\mathop{\mathrm{Spec}}}
\def\Hom{\mathop{\mathrm{Hom}}\nolimits}
\def\Ext{\mathop{\mathrm{Ext}}\nolimits}
\def\SheafHom{\mathop{\mathcal{H}\!\mathit{om}}\nolimits}
\def\SheafExt{\mathop{\mathcal{E}\!\mathit{xt}}\nolimits}
\def\Sch{\mathit{Sch}}
\def\Mor{\mathop{\mathrm{Mor}}\nolimits}
\def\Ob{\mathop{\mathrm{Ob}}\nolimits}
\def\Sh{\mathop{\mathit{Sh}}\nolimits}
\def\NL{\mathop{N\!L}\nolimits}
\def\CH{\mathop{\mathrm{CH}}\nolimits}
\def\proetale{{pro\text{-}\acute{e}tale}}
\def\etale{{\acute{e}tale}}
\def\QCoh{\mathit{QCoh}}
\def\Ker{\mathop{\mathrm{Ker}}}
\def\Im{\mathop{\mathrm{Im}}}
\def\Coker{\mathop{\mathrm{Coker}}}
\def\Coim{\mathop{\mathrm{Coim}}}

% Boxtimes
%
\DeclareMathSymbol{\boxtimes}{\mathbin}{AMSa}{"02}

%
% Macros for moduli stacks/spaces
%
\def\QCohstack{\mathcal{QC}\!\mathit{oh}}
\def\Cohstack{\mathcal{C}\!\mathit{oh}}
\def\Spacesstack{\mathcal{S}\!\mathit{paces}}
\def\Quotfunctor{\mathrm{Quot}}
\def\Hilbfunctor{\mathrm{Hilb}}
\def\Curvesstack{\mathcal{C}\!\mathit{urves}}
\def\Polarizedstack{\mathcal{P}\!\mathit{olarized}}
\def\Complexesstack{\mathcal{C}\!\mathit{omplexes}}
% \Pic is the operator that assigns to X its picard group, usage \Pic(X)
% \Picardstack_{X/B} denotes the Picard stack of X over B
% \Picardfunctor_{X/B} denotes the Picard functor of X over B
\def\Pic{\mathop{\mathrm{Pic}}\nolimits}
\def\Picardstack{\mathcal{P}\!\mathit{ic}}
\def\Picardfunctor{\mathrm{Pic}}
\def\Deformationcategory{\mathcal{D}\!\mathit{ef}}
\usepackage{cleveref}
\crefname{originalproposition}{original proposition}{original propositions}
\Crefname{originalproposition}{Original proposition}{Original propositions}
\makeatletter
\newlabel{proposition-connected-component}{{0B7R}{source}{}{sourcepointer-1}{}}
\newlabel{proposition-connected-component@cref}{{[originalproposition][0][]0B7R}{[1][0][]source}{}{}{}}
\makeatother
\crefname{originallemma}{original lemma}{original lemmas}
\Crefname{originallemma}{Original lemma}{Original lemmas}
\makeatletter
\newlabel{limits-lemma-topology-limit}{{01YY}{source}{}{sourcepointer-2}{}}
\newlabel{limits-lemma-topology-limit@cref}{{[originallemma][0][]01YY}{[1][0][]source}{}{}{}}
\makeatother
\crefname{originaldefinition}{original definition}{original definitions}
\Crefname{originaldefinition}{Original definition}{Original definitions}
\makeatletter
\newlabel{topology-definition-profinite}{{08ZX}{source}{}{sourcepointer-3}{}}
\newlabel{topology-definition-profinite@cref}{{[originaldefinition][0][]08ZX}{[1][0][]source}{}{}{}}
\makeatother
\crefname{originallemma}{original lemma}{original lemmas}
\Crefname{originallemma}{Original lemma}{Original lemmas}
\makeatletter
\newlabel{lemma-group-scheme-over-field-translate-open}{{0B7N}{source}{}{sourcepointer-4}{}}
\newlabel{lemma-group-scheme-over-field-translate-open@cref}{{[originallemma][0][]0B7N}{[1][0][]source}{}{}{}}
\makeatother
\crefname{originallemma}{original lemma}{original lemmas}
\Crefname{originallemma}{Original lemma}{Original lemmas}
\makeatletter
\newlabel{topology-lemma-spectral-pi0}{{0906}{source}{}{sourcepointer-5}{}}
\newlabel{topology-lemma-spectral-pi0@cref}{{[originallemma][0][]0906}{[1][0][]source}{}{}{}}
\makeatother
\crefname{originallemma}{original lemma}{original lemmas}
\Crefname{originallemma}{Original lemma}{Original lemmas}
\makeatletter
\newlabel{properties-lemma-quasi-compact-quasi-separated-spectral}{{094L}{source}{}{sourcepointer-6}{}}
\newlabel{properties-lemma-quasi-compact-quasi-separated-spectral@cref}{{[originallemma][0][]094L}{[1][0][]source}{}{}{}}
\makeatother
\crefname{originallemma}{original lemma}{original lemmas}
\Crefname{originallemma}{Original lemma}{Original lemmas}
\makeatletter
\newlabel{topology-lemma-closed-map}{{08YD}{source}{}{sourcepointer-7}{}}
\newlabel{topology-lemma-closed-map@cref}{{[originallemma][0][]08YD}{[1][0][]source}{}{}{}}
\makeatother
\crefname{originallemma}{original lemma}{original lemmas}
\Crefname{originallemma}{Original lemma}{Original lemmas}
\makeatletter
\newlabel{topology-lemma-profinite-refine-open-covering}{{08ZZ}{source}{}{sourcepointer-8}{}}
\newlabel{topology-lemma-profinite-refine-open-covering@cref}{{[originallemma][0][]08ZZ}{[1][0][]source}{}{}{}}
\makeatother
\crefname{originallemma}{original lemma}{original lemmas}
\Crefname{originallemma}{Original lemma}{Original lemmas}
\makeatletter
\newlabel{lemma-profinite-product-over-field}{{0B7T}{source}{}{sourcepointer-9}{}}
\newlabel{lemma-profinite-product-over-field@cref}{{[originallemma][0][]0B7T}{[1][0][]source}{}{}{}}
\makeatother
\title{The image of a pro-affine scheme in a Jacobson group scheme with rational closed points lands in an affine open}
\author{The Stacks Project Authors}
\date{stacks-project:groupoids.tex:local-snapshot-20261012; GNU FDL 1.2 or later}
\begin{document}
\maketitle
\setcounter{subsection}{1}
\noindent\textit{Editorial note.} All delivered mathematical text is the author's own, copied from the pinned Stacks Project chapter file groupoids.tex; no statement or proof step was rewritten, reordered or removed. The counted span is the author environment and its own proof environment. References inside the span are handed to the assembler as navigation stubs that carry the canonical tag of the cited tag only; whether each of them is load-bearing is recorded on the stub itself.\par\medskip
\noindent\textit{Allowed standard prerequisites (editorial).}\begin{itemize}
\item Proposition proposition-connected-component (canonical tag 0B7R): the connected component G\textasciicircum{}0 is a closed subgroup scheme and components are translates. cited at proof lines 1318-1319. Supplies G\textasciicircum{}0 for the component-wise reduction.
\item Lemma limits-lemma-topology-limit (canonical tag 01YY), Definition topology-definition-profinite (08ZX), Lemma topology-lemma-spectral-pi0 (0906) with properties-lemma-quasi-compact-quasi-separated-spectral (094L). cited at proof lines 1322-1337. T and pi\_0(W) are profinite, which drives the finite clopen decomposition.
\item Lemma topology-lemma-closed-map (canonical tag 08YD) and topology-lemma-profinite-refine-open-covering (08ZZ). cited at proof lines 1340-1349. Closedness of psi produces the neighbourhoods; the refinement makes the cover finite and disjoint.
\item Lemma lemma-group-scheme-over-field-translate-open (canonical tag 0B7N). cited at proof lines 1327-1329. Makes W invariant under left translation by G\textasciicircum{}0.
\item Lemma lemma-profinite-product-over-field (canonical tag 0B7T). cited at proof lines 1382-1386. Identifies |U x\_k T| = |U| x |T| so the image of the open immersion is a finite union of boxes.
\end{itemize}\medskip
\ifdefined\hypertarget\hypertarget{sourcepointer-1}{}\fi
\noindent\textit{Source reference.} Original proposition 0B7R: sources/stacks-groupoids/pinned-original-source.tex, source lines 1230--1230. Cited at proof lines 1318-1319 to supply the closed subgroup scheme G\textasciicircum{}0 used in the two reduction paragraphs; load-bearing.\par\medskip
\ifdefined\hypertarget\hypertarget{sourcepointer-2}{}\fi
\noindent\textit{Source reference.} Original lemma 01YY: sources/stacks-chow/tags.txt, source lines 2573--2573. Cited at proof line 1322 to see T is a directed inverse limit of finite topological spaces; load-bearing for the profinite observation.\par\medskip
\ifdefined\hypertarget\hypertarget{sourcepointer-3}{}\fi
\noindent\textit{Source reference.} Original definition 08ZX: sources/stacks-chow/tags.txt, source lines 11390--11390. Cited at proof lines 1322-1324 to name the topology of T profinite; framing of the topological argument.\par\medskip
\ifdefined\hypertarget\hypertarget{sourcepointer-4}{}\fi
\noindent\textit{Source reference.} Original lemma 0B7N: sources/stacks-groupoids/pinned-original-source.tex, source lines 965--965. Cited at proof lines 1327-1329 to make W invariant under left translation by G\textasciicircum{}0; load-bearing.\par\medskip
\ifdefined\hypertarget\hypertarget{sourcepointer-5}{}\fi
\noindent\textit{Source reference.} Original lemma 0906: sources/stacks-chow/tags.txt, source lines 11399--11399. Cited at proof lines 1334-1337 (with 094L) to see pi\_0(W) is profinite; load-bearing for the finite clopen refinement.\par\medskip
\ifdefined\hypertarget\hypertarget{sourcepointer-6}{}\fi
\noindent\textit{Source reference.} Original lemma 094L: sources/stacks-chow/tags.txt, source lines 11554--11554. Cited at proof lines 1334-1337 (with 0906) to see pi\_0(W) is profinite; load-bearing for the finite clopen refinement.\par\medskip
\ifdefined\hypertarget\hypertarget{sourcepointer-7}{}\fi
\noindent\textit{Source reference.} Original lemma 08YD: sources/stacks-chow/tags.txt, source lines 11336--11336. Cited at proof lines 1340-1342 to see psi is closed, whence the neighbourhoods V\_xi exist; load-bearing.\par\medskip
\ifdefined\hypertarget\hypertarget{sourcepointer-8}{}\fi
\noindent\textit{Source reference.} Original lemma 08ZZ: sources/stacks-chow/tags.txt, source lines 11392--11392. Cited at proof lines 1346-1349 to replace the covering by V\_1, ..., V\_n; load-bearing.\par\medskip
\ifdefined\hypertarget\hypertarget{sourcepointer-9}{}\fi
\noindent\textit{Source reference.} Original lemma 0B7T: sources/stacks-groupoids/pinned-original-source.tex, source lines 1269--1269. Cited at proof lines 1382-1386 to read |U x\_k T| = |U| x |T|, whence the image open immersion is a finite union of boxes; load-bearing.\par\medskip
\begin{lemma}
\label{lemma-compact-set-in-affine}
Let $k$ be an algebraically closed field. Let $G$ be a group scheme over $k$.
Assume that $G$ is Jacobson and that all closed points are $k$-rational.
Let $T = \Spec(A)$ where $A$ is a directed colimit of algebras which
are finite products of copies of $k$. For any morphism $f : T \to G$
there exists an affine open $U \subset G$ containing $f(T)$.
\end{lemma}

\begin{proof}
Let $G^0 \subset G$ be the closed subgroup scheme found in
Proposition \ref{proposition-connected-component}. The first two paragraphs
serve to reduce to the case $G = G^0$.

\medskip\noindent
Observe that $T$ is a directed inverse limit of finite topological spaces
(Limits, Lemma \ref{limits-lemma-topology-limit}), hence profinite as a
topological space (Topology, Definition \ref{topology-definition-profinite}).
Let $W \subset G$ be a quasi-compact open containing the image of $T \to G$.
After replacing $W$ by the image of $G^0 \times W \to G \times G \to G$ we may
assume that $W$ is invariant under the action of left translation by $G^0$, see
Lemma \ref{lemma-group-scheme-over-field-translate-open}.
Consider the composition
$$
\psi = \pi \circ f : T \xrightarrow{f} W \xrightarrow{\pi} \pi_0(W)
$$
The space $\pi_0(W)$ is profinite
(Topology, Lemma \ref{topology-lemma-spectral-pi0} and
Properties, Lemma
\ref{properties-lemma-quasi-compact-quasi-separated-spectral}).
Let $F_\xi \subset T$ be the fibre of $T \to \pi_0(W)$ over $\xi \in \pi_0(W)$.
Suppose that for every $\xi \in \pi_0(W)$ we have found affine open
$U_\xi \subset W$ such that $f(F_\xi) \subset U_\xi$. As $\psi$ is closed by
Topology, Lemma \ref{topology-lemma-closed-map},
$\psi(T \setminus f^{-1}(U_\xi))$ is a closed subset of $\pi_0(W)$
that does not contain $\xi$. Hence its complement $V_\xi$ is an
open neighborhood of $\xi$ and $\psi^{-1}(V_\xi) \subseteq f^{-1}(U_\xi)$.
The open covering of $\pi_0(W)$ by the $V_\xi$ may be refined by
a finite covering $\pi_0(W) = V_1 \amalg \ldots \amalg V_n$
of pairwise disjoint open and closed subsets (Topology,
Lemma \ref{topology-lemma-profinite-refine-open-covering}).
For every $i = 1, \ldots, n$ pick $\xi_i \in \pi_0(W)$
such that $V_i \subset V_{\xi_i}$. As $V_i$ is closed and
$U_{\xi_i}$ affine, $U_{\xi_i} \cap \pi^{-1}(V_i)$ is affine as well.
Hence replacing $U_{\xi_i}$ by $U_{\xi_i} \cap \pi^{-1}(V_i)$ we
have $\psi^{-1}(V_i) = f^{-1}(U_{\xi_i})$. As the $V_i$'s cover
$\pi_0(W)$, the $f^{-1}(U_{\xi_i})$'s cover $T$ and so
$f(T) \subseteq U = \bigcup_i U_{\xi_i}$. Then $U$ is the disjoint
union of finitely many affine opens and hence affine as desired.

\medskip\noindent
Let $Z$ be a connected component of $G$ which meets $f(T)$. Then $Z$
has a $k$-rational point $z$ (because all residue fields of the scheme $T$
are isomorphic to $k$). Hence $Z = G^0 z$. By our choice of $W$, we see
that $Z \subset W$. The argument in the preceding paragraph reduces us to
the problem of finding an affine open neighbourhood of $f(T) \cap Z$ in $W$.
After translation by a rational point we may assume that $Z = G^0$
(details omitted). Observe that the scheme theoretic inverse image
$T' = f^{-1}(G^0) \subset T$ is a closed subscheme, which has the same type.
After replacing $T$ by $T'$ we may assume that $f(T) \subset G^0$.
Choose an affine open neighbourhood $U \subset G$
of $e \in G$, so that in particular $U \cap G^0$ is nonempty. We will show
there exists a $g \in G^0(k)$ such that $f(T) \subset g^{-1}U$.
This will finish the proof as $g^{-1}U \subset W$ by the left
$G^0$-invariance of $W$.

\medskip\noindent
The arguments in the preceding two paragraphs allow us to pass to $G^0$
and reduce the problem to the following:
Assume $G$ is irreducible and $U \subset G$ an affine
open neighbourhood of $e$. Show that $f(T) \subset g^{-1}U$
for some $g \in G(k)$. Consider the morphism
$$
U \times_k T \longrightarrow G \times_k T,\quad
(t, u) \longrightarrow (uf(t)^{-1}, t)
$$
which is an open immersion (because the extension of this morphism to
$G \times_k T \to G \times_k T$ is an isomorphism).
By our assumption on $T$ we see that we have $|U \times_k T| = |U| \times |T|$
and similarly for $G \times_k T$, see
Lemma \ref{lemma-profinite-product-over-field}.
Hence the image of the displayed open immersion is a finite union
of boxes $\bigcup_{i = 1, \ldots, n} U_i \times V_i$ with
$V_i \subset T$ and $U_i \subset G$ quasi-compact open. This means that
the possible opens $Uf(t)^{-1}$, $t \in T$ are finite in number, say
$Uf(t_1)^{-1}, \ldots, Uf(t_r)^{-1}$. Since $G$ is irreducible the
intersection
$$
Uf(t_1)^{-1} \cap \ldots \cap Uf(t_r)^{-1}
$$
is nonempty and since $G$ is Jacobson with closed points $k$-rational,
we can choose a $k$-valued point $g \in G(k)$ of this intersection.
Then we see that $g \in Uf(t)^{-1}$ for all $t \in T$ which
means that $f(t) \in g^{-1}U$ as desired.
\end{proof}

\end{document}
